\documentclass{article}

\usepackage{amsmath, amsthm, amssymb, amsfonts}
\usepackage{graphicx}
\usepackage{setspace}
\usepackage{geometry}
\usepackage[utf8]{inputenc}
\usepackage[english]{babel}
\usepackage{bbold}

\newcommand{\HRule}[1]{\rule{\linewidth}{#1}}

\setstretch{1.2}
\geometry{
    textheight=9in,
    textwidth=5.5in,
    top=1in,
    headheight=12pt,
    headsep=25pt,
    footskip=30pt
}

% ------------------------------------------------------------------------------

\begin{document}

\title{ \normalsize \textsc{}
		\\ [2.0cm]
		\HRule{1.5pt} \\
		\LARGE \textbf{\uppercase{Quantum Molecular Hamiltonian}
		\HRule{2.0pt} \\ [0.6cm] \LARGE{From Lab-Frame Cartesian to Total-CoM Coordinates} \vspace*{10\baselineskip}}
		}
\date{}
\author{\textbf{Chih-En Shen}}

\maketitle
\newpage

\tableofcontents
\newpage

% ------------------------------------------------------------------------------

\section{Quantum Hamiltonian in Cartesian coordinates}

As in the classical treatment, we have $N_{\mathrm{nucl}}$ nuclei labeled by $j$ with masses $M_j$, charges $Z_j e$, and positions $\mathbf{R}_j$; and $N_{\mathrm{elec}}$ electrons labeled by $i$ with mass $m_e$, charge $-e$, and positions $\mathbf{r}_i$. All vectors live in a space-fixed (SF) lab frame. The state is a wavefunction $\Psi(\mathbf{R}_1,\dots,\mathbf{R}_{N_{\mathrm{nucl}}},\mathbf{r}_1,\dots,\mathbf{r}_{N_{\mathrm{elec}}})$, and canonical quantization in Cartesian coordinates is the substitution $\mathbf{P}_j\to\hat{\mathbf{P}}_j=-i\hbar\nabla_{\mathbf{R}_j}$, $\mathbf{p}_i\to\hat{\mathbf{p}}_i=-i\hbar\nabla_{\mathbf{r}_i}$, giving

\begin{equation}
\hat H \;=\; \underbrace{-\sum_{j=1}^{N_{\mathrm{nucl}}}\frac{\hbar^2}{2M_j}\nabla^2_{\mathbf{R}_j}}_{\hat T_N} \;\underbrace{-\sum_{i=1}^{N_{\mathrm{elec}}}\frac{\hbar^2}{2m_e}\nabla^2_{\mathbf{r}_i}}_{\hat T_e} \;+\; V, \tag{1.1}
\end{equation}
where $\nabla^2_{\mathbf{R}_j}=\sum_{\alpha=x,y,z}\partial^2/\partial R_{j\alpha}^2$ and, writing $R_{ij}=|\mathbf{R}_i-\mathbf{R}_j|$ etc.,
\begin{equation}
V = \frac{1}{4\pi\epsilon_0}\left(\frac12\sum_{i\ne j}\frac{Z_i Z_j e^2}{R_{ij}} - \sum_{i,j}\frac{Z_j e^2}{|\mathbf{R}_j-\mathbf{r}_i|} + \frac12\sum_{i\ne j}\frac{e^2}{r_{ij}}\right). \tag{1.2}
\end{equation}
The canonical commutators are $[\hat R_{j\alpha},\hat P_{k\beta}]=i\hbar\,\delta_{jk}\delta_{\alpha\beta}$, $[\hat r_{i\alpha},\hat p_{l\beta}]=i\hbar\,\delta_{il}\delta_{\alpha\beta}$, and the inner product is the flat Cartesian one,
\begin{equation}
\langle\Phi|\Psi\rangle=\int \Phi^{*}\,\Psi\;\prod_j d^{3}R_j\;\prod_i d^{3}r_i , \tag{1.3}
\end{equation}
with respect to which each $-i\hbar\,\partial$ is Hermitian by integration by parts (wavefunctions vanishing at infinity). No ordering ambiguity arises in (1.1): the masses are constants, so the kinetic operators are unambiguous.

The task of this note is purely computational: re-express the \emph{derivatives} in (1.1) in total-CoM coordinates using the chain rule, and watch the CoM motion separate.

\section{Total-CoM coordinates and the transformed wavefunction}

\subsection{New coordinates}

Exactly as in the classical treatment, define the \emph{total} (nuclei $+$ electrons) CoM and refer every particle to it:
\begin{equation}
\mathbf{R}_{\mathrm{cm}} = \frac{1}{M_{\mathrm{tot}}}\Bigl(\sum_j M_j\mathbf{R}_j + m_e\sum_i\mathbf{r}_i\Bigr),\qquad M_{\mathrm{tot}} = \underbrace{\sum_j M_j}_{\equiv\,M_N} +\, N_{\mathrm{elec}}\,m_e, \tag{2.1}
\end{equation}
\begin{equation}
\mathbf{R}'_j = \mathbf{R}_j - \mathbf{R}_{\mathrm{cm}},\qquad \mathbf{r}'_i = \mathbf{r}_i - \mathbf{R}_{\mathrm{cm}}. \tag{2.2}
\end{equation}
Multiplying each equation in (2.2) by its mass and summing, the definition (2.1) forces the linear relation
\begin{equation}
\sum_j M_j\mathbf{R}'_j + m_e\sum_i\mathbf{r}'_i = 0. \tag{2.3}
\end{equation}

\subsection{The redundancy and how the wavefunction handles it}

Count coordinates. The old set has $3(N_{\mathrm{nucl}}+N_{\mathrm{elec}})$ components; the new set $(\mathbf{R}_{\mathrm{cm}},\{\mathbf{R}'_j\},\{\mathbf{r}'_i\})$ has $3(N_{\mathrm{nucl}}+N_{\mathrm{elec}}+1)$, three too many, and (2.3) is precisely the compensating relation. We express the wavefunction in the new variables,
\begin{equation}
\Psi(\mathbf{R}_j,\mathbf{r}_i)\;=\;\psi\bigl(\mathbf{R}_{\mathrm{cm}},\{\mathbf{R}'_j\},\{\mathbf{r}'_i\}\bigr), \tag{2.4}
\end{equation}
but because the new arguments are constrained by (2.3), equation (2.4) does not define $\psi$ \emph{off} the constraint set: $\psi$ is free there, and every choice represents the same physical $\Psi$. Before we can differentiate $\psi$ with respect to its arguments \emph{as if they were independent} --- which is what the chain rule requires --- we must complete its definition. Introduce the total internal-gradient operator
\begin{equation}
\mathbf{D}\;\equiv\;\sum_k\nabla_{\mathbf{R}'_k}\;+\;\sum_l\nabla_{\mathbf{r}'_l}, \tag{2.5}
\end{equation}
which generates an \emph{equal shift} of all relative vectors, $\mathbf{R}'_k\to\mathbf{R}'_k+\boldsymbol{\varepsilon}$, $\mathbf{r}'_l\to\mathbf{r}'_l+\boldsymbol{\varepsilon}$, and complete the definition of $\psi$ by the auxiliary condition
\begin{equation}
\boxed{\;\mathbf{D}\,\psi=0.\;} \tag{2.6}
\end{equation}
This is consistent: under the equal shift the constraint function $\sum_j M_j\mathbf{R}'_j+m_e\sum_i\mathbf{r}'_i$ changes by $M_{\mathrm{tot}}\boldsymbol{\varepsilon}\neq0$, so the shift moves \emph{transversally off} the constraint set --- condition (2.6) only prescribes how $\psi$ continues off the set and places no restriction on its physical values on it. Equation (2.6) is the quantum image of the classical dual constraint $\sum_k\mathbf{P}'_k+\sum_l\mathbf{p}'_l=0$: acting with $-i\hbar$, it reads $\bigl(\sum_k\hat{\mathbf{P}}'_k+\sum_l\hat{\mathbf{p}}'_l\bigr)\psi=0$.

\section{Chain rule: first derivatives}

\subsection{Jacobian of the coordinate change}

Every new coordinate in (2.1)--(2.2) is a \emph{linear} function of the old ones, so all partial derivatives of new with respect to old are constants, diagonal in the Cartesian index. Differentiating (2.1)--(2.2) with respect to a nuclear position,
\begin{equation}
\frac{\partial\mathbf{R}_{\mathrm{cm}}}{\partial\mathbf{R}_j} = \frac{M_j}{M_{\mathrm{tot}}}\mathbb{1},\qquad \frac{\partial\mathbf{R}'_k}{\partial\mathbf{R}_j} = \Bigl(\delta_{jk}-\frac{M_j}{M_{\mathrm{tot}}}\Bigr)\mathbb{1},\qquad \frac{\partial\mathbf{r}'_l}{\partial\mathbf{R}_j} = -\frac{M_j}{M_{\mathrm{tot}}}\mathbb{1}, \tag{3.1a}
\end{equation}
and with respect to an electron position,
\begin{equation}
\frac{\partial\mathbf{R}_{\mathrm{cm}}}{\partial\mathbf{r}_i} = \frac{m_e}{M_{\mathrm{tot}}}\mathbb{1},\qquad \frac{\partial\mathbf{R}'_k}{\partial\mathbf{r}_i} = -\frac{m_e}{M_{\mathrm{tot}}}\mathbb{1},\qquad \frac{\partial\mathbf{r}'_l}{\partial\mathbf{r}_i} = \Bigl(\delta_{il}-\frac{m_e}{M_{\mathrm{tot}}}\Bigr)\mathbb{1}. \tag{3.1b}
\end{equation}

\subsection{Chain rule, written out once in components}

For one Cartesian component $\alpha$ of one nuclear position, the chain rule applied to (2.4) reads
\begin{equation}
\frac{\partial\Psi}{\partial R_{j\alpha}}
=\sum_\beta\frac{\partial R_{\mathrm{cm},\beta}}{\partial R_{j\alpha}}\frac{\partial\psi}{\partial R_{\mathrm{cm},\beta}}
+\sum_k\sum_\beta\frac{\partial R'_{k\beta}}{\partial R_{j\alpha}}\frac{\partial\psi}{\partial R'_{k\beta}}
+\sum_l\sum_\beta\frac{\partial r'_{l\beta}}{\partial R_{j\alpha}}\frac{\partial\psi}{\partial r'_{l\beta}}. \tag{3.2}
\end{equation}
All the Jacobian blocks (3.1) are proportional to $\mathbb{1}$, i.e.\ to $\delta_{\alpha\beta}$, so the $\beta$-sums collapse and the three vector components assemble into
\begin{equation}
\nabla_{\mathbf{R}_j}\Psi=\Bigl[\frac{M_j}{M_{\mathrm{tot}}}\nabla_{\mathbf{R}_{\mathrm{cm}}}+\nabla_{\mathbf{R}'_j}-\frac{M_j}{M_{\mathrm{tot}}}\,\mathbf{D}\Bigr]\psi, \tag{3.3}
\end{equation}
where the $-({M_j}/{M_{\mathrm{tot}}})\mathbf{D}$ collects the $-M_j/M_{\mathrm{tot}}$ entries that appear in \emph{every} relative coordinate in (3.1a). The same steps for an electron position give
\begin{equation}
\nabla_{\mathbf{r}_i}\Psi=\Bigl[\frac{m_e}{M_{\mathrm{tot}}}\nabla_{\mathbf{R}_{\mathrm{cm}}}+\nabla_{\mathbf{r}'_i}-\frac{m_e}{M_{\mathrm{tot}}}\,\mathbf{D}\Bigr]\psi. \tag{3.4}
\end{equation}
Now use the auxiliary condition (2.6). All the operators in (3.3)--(3.4) have \emph{constant} coefficients, so $\mathbf{D}$ commutes with each of them; hence if $\mathbf{D}\psi=0$, then $\mathbf{D}$ annihilates every constant-coefficient derivative of $\psi$ as well, and the $\mathbf{D}$-terms may be dropped \emph{inside any subsequent derivative}:
\begin{equation}
\boxed{\;\nabla_{\mathbf{R}_j}\Psi=\Bigl[\frac{M_j}{M_{\mathrm{tot}}}\nabla_{\mathbf{R}_{\mathrm{cm}}}+\nabla_{\mathbf{R}'_j}\Bigr]\psi,\qquad
\nabla_{\mathbf{r}_i}\Psi=\Bigl[\frac{m_e}{M_{\mathrm{tot}}}\nabla_{\mathbf{R}_{\mathrm{cm}}}+\nabla_{\mathbf{r}'_i}\Bigr]\psi.\;} \tag{3.5}
\end{equation}
Multiplying (3.5) by $-i\hbar$ shows it is precisely the operator image of the classical momentum relation: $\hat{\mathbf{P}}_j=({M_j}/{M_{\mathrm{tot}}})\hat{\mathbf{P}}_{\mathrm{cm}}+\hat{\mathbf{P}}'_j$ and $\hat{\mathbf{p}}_i=({m_e}/{M_{\mathrm{tot}}})\hat{\mathbf{P}}_{\mathrm{cm}}+\hat{\mathbf{p}}'_i$, with $\hat{\mathbf{P}}_{\mathrm{cm}}=-i\hbar\nabla_{\mathbf{R}_{\mathrm{cm}}}$ etc.

\subsection{Cross-check: the total momentum}

Sum (3.3) over all nuclei and (3.4) over all electrons \emph{before} imposing (2.6). The CoM coefficients add to $(M_N+N_{\mathrm{elec}}m_e)/M_{\mathrm{tot}}=1$; the plain relative gradients add to $\mathbf{D}$; and the $\mathbf{D}$-corrections add to $-[(M_N+N_{\mathrm{elec}}m_e)/M_{\mathrm{tot}}]\mathbf{D}=-\mathbf{D}$. The last two cancel \emph{identically} --- no auxiliary condition needed:
\begin{equation}
\Bigl(\sum_j\nabla_{\mathbf{R}_j}+\sum_i\nabla_{\mathbf{r}_i}\Bigr)\Psi=\nabla_{\mathbf{R}_{\mathrm{cm}}}\,\psi
\quad\Longleftrightarrow\quad
\sum_j\hat{\mathbf{P}}_j+\sum_i\hat{\mathbf{p}}_i=\hat{\mathbf{P}}_{\mathrm{cm}}. \tag{3.6}
\end{equation}
The total momentum operator is exactly the CoM gradient, independent of how $\psi$ was extended off the constraint set.

\section{Second derivatives and the transformed Hamiltonian}

\subsection{Laplacians}

Relation (3.5) expresses $\nabla_{\mathbf{R}_j}$, acting on any function written in the new coordinates, as a fixed constant-coefficient combination of new-coordinate derivatives. Applying it twice --- legitimate precisely because the coefficients are constants, and mixed partial derivatives of smooth functions commute --- and expanding the square,
\begin{align}
\nabla^2_{\mathbf{R}_j}\Psi
&=\Bigl[\frac{M_j}{M_{\mathrm{tot}}}\nabla_{\mathbf{R}_{\mathrm{cm}}}+\nabla_{\mathbf{R}'_j}\Bigr]\!\cdot\!\Bigl[\frac{M_j}{M_{\mathrm{tot}}}\nabla_{\mathbf{R}_{\mathrm{cm}}}+\nabla_{\mathbf{R}'_j}\Bigr]\psi \nonumber\\
&=\Bigl[\frac{M_j^2}{M_{\mathrm{tot}}^2}\nabla^2_{\mathbf{R}_{\mathrm{cm}}}
+\frac{2M_j}{M_{\mathrm{tot}}}\,\nabla_{\mathbf{R}_{\mathrm{cm}}}\!\cdot\!\nabla_{\mathbf{R}'_j}
+\nabla^2_{\mathbf{R}'_j}\Bigr]\psi , \tag{4.1}
\end{align}
and identically for the electrons with $M_j\to m_e$, $\mathbf{R}'_j\to\mathbf{r}'_i$. Assemble the nuclear kinetic operator, using $\sum_j M_j=M_N$ in the first term:
\begin{equation}
\hat T_N\Psi=-\sum_j\frac{\hbar^2}{2M_j}\nabla^2_{\mathbf{R}_j}\Psi
=\Bigl[-\frac{M_N}{M_{\mathrm{tot}}}\frac{\hbar^2}{2M_{\mathrm{tot}}}\nabla^2_{\mathbf{R}_{\mathrm{cm}}}
-\frac{\hbar^2}{M_{\mathrm{tot}}}\nabla_{\mathbf{R}_{\mathrm{cm}}}\!\cdot\!\sum_j\nabla_{\mathbf{R}'_j}
-\sum_j\frac{\hbar^2}{2M_j}\nabla^2_{\mathbf{R}'_j}\Bigr]\psi , \tag{4.2}
\end{equation}
\begin{equation}
\hat T_e\Psi=-\sum_i\frac{\hbar^2}{2m_e}\nabla^2_{\mathbf{r}_i}\Psi
=\Bigl[-\frac{N_{\mathrm{elec}}m_e}{M_{\mathrm{tot}}}\frac{\hbar^2}{2M_{\mathrm{tot}}}\nabla^2_{\mathbf{R}_{\mathrm{cm}}}
-\frac{\hbar^2}{M_{\mathrm{tot}}}\nabla_{\mathbf{R}_{\mathrm{cm}}}\!\cdot\!\sum_i\nabla_{\mathbf{r}'_i}
-\sum_i\frac{\hbar^2}{2m_e}\nabla^2_{\mathbf{r}'_i}\Bigr]\psi . \tag{4.3}
\end{equation}
Adding (4.2) and (4.3): the two mass fractions in front of $\nabla^2_{\mathbf{R}_{\mathrm{cm}}}$ sum to $(M_N+N_{\mathrm{elec}}m_e)/M_{\mathrm{tot}}=1$, and the two cross terms combine into the total internal gradient (2.5),
\begin{equation}
-\frac{\hbar^2}{M_{\mathrm{tot}}}\,\nabla_{\mathbf{R}_{\mathrm{cm}}}\!\cdot\!\underbrace{\Bigl(\sum_j\nabla_{\mathbf{R}'_j}+\sum_i\nabla_{\mathbf{r}'_i}\Bigr)}_{=\,\mathbf{D}}\psi \;=\; 0
\qquad\text{by the auxiliary condition (2.6),} \tag{4.4}
\end{equation}
where dropping $\mathbf{D}$ under the extra $\nabla_{\mathbf{R}_{\mathrm{cm}}}$ is again justified because $\mathbf{D}$ commutes with constant-coefficient derivatives. No CoM--internal cross-coupling survives.

\subsection{The potential}

$V$ is a multiplication operator, so no derivatives act on it during the coordinate change; we only rewrite its arguments. Subtracting (2.2) pairwise, $\mathbf{R}'_j-\mathbf{R}'_k = \mathbf{R}_j-\mathbf{R}_k$, $\mathbf{R}'_j-\mathbf{r}'_i = \mathbf{R}_j-\mathbf{r}_i$, $\mathbf{r}'_i-\mathbf{r}'_l = \mathbf{r}_i-\mathbf{r}_l$: every interparticle separation in (1.2) equals the corresponding primed separation, so $V$ is unchanged in form,
\begin{equation}
V=V(\mathbf{R}'_j,\mathbf{r}'_i). \tag{4.5}
\end{equation}

\subsection{Transformed Hamiltonian}

Collecting (4.2)--(4.5), the Schr\"odinger operator in total-CoM coordinates is
\begin{equation}
\boxed{\;\hat H\,\psi=\Bigl[-\frac{\hbar^2}{2M_{\mathrm{tot}}}\nabla^2_{\mathbf{R}_{\mathrm{cm}}}
-\sum_j\frac{\hbar^2}{2M_j}\nabla^2_{\mathbf{R}'_j}
-\sum_i\frac{\hbar^2}{2m_e}\nabla^2_{\mathbf{r}'_i}
+V(\mathbf{R}'_j,\mathbf{r}'_i)\Bigr]\psi,\qquad \mathbf{D}\psi=0.\;} \tag{4.6}
\end{equation}
Exactly as in the classical result: the CoM kinetic term carries the correct total mass $M_{\mathrm{tot}}$, the internal kinetic energy is \emph{diagonal} --- no mass-polarization term --- and no cross-coupling appears, for \emph{any} CoM motion. The price is the constraint pair, (2.3) on the coordinates and (2.6) on the wavefunction.

\subsection{Consistency: $\hat H$ preserves the auxiliary condition}

Condition (2.6) is only meaningful if $\hat H$ maps the space $\{\psi:\mathbf{D}\psi=0\}$ into itself, i.e.\ if $[\mathbf{D},\hat H]=0$. The kinetic terms in (4.6) are constant-coefficient derivatives, which commute with $\mathbf{D}$. For the potential, $\mathbf{D}$ is a first-order operator, so $[\mathbf{D},V]\psi=(\mathbf{D}V)\psi$, and
\begin{equation}
\mathbf{n}\cdot\mathbf{D}V=\mathbf{n}\cdot\Bigl(\sum_k\nabla_{\mathbf{R}'_k}V+\sum_l\nabla_{\mathbf{r}'_l}V\Bigr)=\frac{d}{d\varepsilon}\,V\bigl(\mathbf{R}'_j+\varepsilon\mathbf{n},\,\mathbf{r}'_i+\varepsilon\mathbf{n}\bigr)\Big|_{\varepsilon=0}=0 \tag{4.7}
\end{equation}
for every direction $\mathbf{n}$, because by (4.5) $V$ depends only on \emph{differences} of the primed vectors, which the equal shift leaves untouched; hence $\mathbf{D}V=0$. Hence $[\mathbf{D},\hat H]=0$: a wavefunction satisfying (2.6) keeps satisfying it under the dynamics.

\subsection{Separation of the CoM motion}

$\hat H$ in (4.6) contains $\mathbf{R}_{\mathrm{cm}}$ only through $\nabla^2_{\mathbf{R}_{\mathrm{cm}}}$, so it commutes with $\hat{\mathbf{P}}_{\mathrm{cm}}=-i\hbar\nabla_{\mathbf{R}_{\mathrm{cm}}}$ and the ansatz
\begin{equation}
\psi=e^{\,i\mathbf{K}\cdot\mathbf{R}_{\mathrm{cm}}}\;\chi\bigl(\{\mathbf{R}'_j\},\{\mathbf{r}'_i\}\bigr) \tag{4.8}
\end{equation}
separates exactly: every $\nabla_{\mathbf{R}_{\mathrm{cm}}}$ pulls down $i\mathbf{K}$ and touches nothing else (note $\mathbf{D}$ does not act on $\mathbf{R}_{\mathrm{cm}}$, so $\mathbf{D}\psi=0\Leftrightarrow\mathbf{D}\chi=0$). The eigenvalue problem $\hat H\psi=E\psi$ becomes
\begin{equation}
\boxed{\;\hat H_{\mathrm{int}}\,\chi=\Bigl[-\sum_j\frac{\hbar^2}{2M_j}\nabla^2_{\mathbf{R}'_j}
-\sum_i\frac{\hbar^2}{2m_e}\nabla^2_{\mathbf{r}'_i}
+V(\mathbf{R}'_j,\mathbf{r}'_i)\Bigr]\chi
=\Bigl(E-\frac{\hbar^2K^2}{2M_{\mathrm{tot}}}\Bigr)\chi,\qquad \mathbf{D}\chi=0, \;} \tag{4.9}
\end{equation}
a free particle of mass $M_{\mathrm{tot}}$ carrying the CoM, times an internal problem whose kinetic energy is diagonal.

\subsection{Remark: why no Podolsky machinery is needed here}

The transformation (2.1)--(2.2) is \emph{linear with constant coefficients}. Choosing any independent subset of the new coordinates (say, dropping one relative vector via (2.3)), the old-to-new map is an invertible linear map with a \emph{constant} Jacobian determinant, so the Cartesian volume element (1.3) goes over to the flat new-coordinate volume element times an irrelevant constant. The weight function in the inner product stays constant, $-i\hbar\,\partial$ remains Hermitian by the same integration by parts as in Section 1, and the naive substitution rule survives the coordinate change untouched --- which is why plain chain-rule calculus sufficed here. This is the essential contrast with \emph{curvilinear} coordinates (Euler angles, normal coordinates), where the Jacobian depends on the coordinates, the volume element acquires a nontrivial weight $\sqrt{g}$, and the kinetic operator must instead be built from the Laplace--Beltrami (Podolsky) construction, as carried out in the main text.

\end{document}
